\section{Haar projections and posterior replica information}
\label{sec:haar}

This comparison uses Haar frames rather than the Gaussian side matrices
of Sections~\ref{sec:gaussian} and~\ref{sec:distance-bins}.  We express the
replica measure by successive simplex heights.  The normalization changes
at each insertion because the remaining ambient dimension shrinks.
A projected-distance distribution-function test, valid even with atoms,
recovers a fraction of the height cost.  We choose its exponent below one
as a function of $d$ and keep the resulting entropy remainder explicit.

Proposition~\ref{haar:block} gives that remainder in terms of the state
entropies.  In particular, the block cost is $O(d)$ when all state
alphabets have logarithmic width $o(d^2)$.  We first prove the fixed-prior
height identity and cancellation, then average over state values.  A Gaussian observation
block is converted to its orthonormal row frame only in that final step;
the learner's observations are still the original exact Gaussian pairs.

Write $\sigma=\sigma_d$ and $\log^+x=\max\{0,\log x\}$ for $x>0$.
A $j$-frame is a
$j\times d$ matrix with orthonormal rows; its Haar law is the probability
law invariant under orthogonal changes of ambient coordinates.  Throughout
the section set
\begin{equation}\label{haar:dimensions}
 d\ge64,\qquad k=r=\lfloor d/16\rfloor,\qquad m=r+1,
 \qquad \ell=\lfloor d/4\rfloor.
\end{equation}
Here $k$ is the actual block row count, $r$ is the number of height
insertions, $m=r+1$ is the replica count, and $\ell$ is the side-frame
row count.
At the $i$th insertion, $1\le i\le r$, the two dimensions that
will matter are
\begin{equation}\label{haar:dimension-slack}
 n_i=d-i+1,\qquad q_i=\ell-i+1.
\end{equation}
They satisfy
\begin{equation}\label{haar:dimension-bounds}
 n_i>k+2,\qquad q_i\ge\frac{3d}{16}>k,
 \qquad \frac{q_i}{q_i-k}\le\frac32,
 \qquad \frac{n_i}{q_i}\le\frac{16}{3},\qquad \frac d k\le32.
\end{equation}
Indeed, $q_i\ge\ell-r+1\ge3d/16$ and $k\le d/16$; also
$k\ge d/32$ for $d\ge64$.  The other inequalities follow at once.  In
particular, none of the denominators in \eqref{haar:dimension-bounds}
vanishes.

\subsection{An exact identity for the replica measure}

We begin with the normalization of spherical projection.  For integers
$1\le j<n-2$, put
\[
 \alpha_{n,j}=\frac{\Gamma(n/2)}{\pi^{j/2}\Gamma((n-j)/2)},
 \qquad v_j=\frac{\pi^{j/2}}{\Gamma(j/2+1)}.
\]
Here $v_j$ is the volume of the unit ball in $\mathbb R^j$.  The first
$j$ coordinates of a uniform point of $S^{n-1}$ have density
\begin{equation}\label{haar:sphere-density}
 p_{n,j}(u)=\alpha_{n,j}(1-\lVert u\rVert^2)^{(n-j-2)/2}
            \mathbf 1_{\{\lVert u\rVert<1\}}.
\end{equation}
Conditional on these coordinates, the remaining direction is uniform,
with radius $\sqrt{1-\lVert u\rVert^2}$.  These are the density and
conditional law in Lemma~\ref{lem:projection-fibers}, specialized to
ambient dimension $n$ and a matrix with orthonormal rows.  We keep the
ambient dimension in the notation because it changes at each insertion.

We will also use the elementary estimates
\begin{equation}\label{haar:gamma-bounds}
 \alpha_{n,j}\le n^{j/2},\qquad
 v_j\alpha_{n,j}\le (Cn/j)^{j/2},
\end{equation}
where $C$ is absolute.  For completeness, gamma recursion and
$\Gamma(z+1/2)\le\sqrt z\,\Gamma(z)$ imply
$\Gamma(n/2)/\Gamma((n-j)/2)\le(n/2)^{j/2}$; the half-step inequality
is Cauchy--Schwarz in the gamma integral.  Integrating that integral over
$[j/2,j/2+1]$ gives
$\Gamma(j/2+1)\ge(j/2)^{j/2}e^{-j/2-1}$.  These two bounds prove
\eqref{haar:gamma-bounds}.

Let $P=g\sigma$ be a probability measure with a bounded Borel density
$0\le g\le K<\infty$.  Necessarily $K\ge1$.  Draw an independent Haar
$k$-frame $A$ and $S\sim P$, and write $Y=AS$.  Under $\sigma$, the
density of $Y$ is
\[
 p^0(y)=\alpha_{d,k}(1-\lVert y\rVert^2)^{(d-k-2)/2}
        \mathbf 1_{\{\lVert y\rVert<1\}}.
\]
For $\lVert y\rVert<1$, let $U_{A,y}$ be uniform probability on the
sphere $\{s\in S^{d-1}:As=y\}$.  On $\lVert y\rVert\ge1$, choose any
fixed measurable probability kernel for $U_{A,y}$; it will always be
multiplied by $p^0(y)=0$.  Fix the following density and posterior
versions:
\begin{equation}\label{haar:posterior}
 \begin{split}
 p_A(y)&=p^0(y)\int g\,\mathrm dU_{A,y}
             \quad (\lVert y\rVert<1),\qquad
 p_A(y)=0\quad (\lVert y\rVert\ge1),\\
 P_{A,y}(\mathrm ds)&=
    \frac{g(s)}{\int g\,\mathrm dU_{A,y}}\,U_{A,y}(\mathrm ds)
       \quad\text{when }p_A(y)>0.
 \end{split}
\end{equation}
The posterior on the unused set $\{p_A=0\}$ may be any fixed measurable
probability kernel.  These are measurable versions.  For example, a
fiber direction can be generated by projecting a standard Gaussian onto
$\ker A$ and normalizing it.  If $u$ is that unit direction, the fiber
point is $A^{\mathsf T}y+\sqrt{1-\lVert y\rVert^2}\,u$.  A fixed
coordinate convention handles the null normalization failures.  This
construction
also makes $p_A(y)$ jointly measurable in $(A,y)$.  We have
$p_A\le Kp^0$ and, for every $A$, $p_A$ is a probability density.

For a tuple $\boldsymbol s=(s_0,\ldots,s_r)$, define
\[
 D_i=\operatorname{span}\{s_1-s_0,\ldots,s_i-s_0\},\qquad D_0=\{0\},
 \qquad
 \delta_i=\operatorname{dist}(s_i-s_0,D_{i-1}).
\]
Thus $\delta_i$ is the height added by the $i$th point.  Write
$\mathrm dA_D$ for Haar probability on the $k$-frames in $D^\perp$.
All subspaces arising here have $\dim D^\perp\ge d-r>k$.  A measurable
version of this frame kernel is obtained by projecting independent
standard Gaussians onto $D^\perp$ and applying a fixed Gram--Schmidt
rule.  We use $\mathrm dA$ when $D=\{0\}$.

\begin{lemma}[Haar simplex-height identity]\label{haar:coarea}
For the dimensions in \eqref{haar:dimensions} and every probability
$P=g\sigma$ with $0\le g\le K<\infty$, the following finite measures on
the frame and the whole tuple are equal:
\begin{equation}\label{haar:coarea-identity}
 \mathrm dA\int_{\mathbb R^k}p_A(y)^m\,\mathrm dy
       \prod_{i=0}^r P_{A,y}(\mathrm ds_i)
 =\left(\prod_{i=0}^rP(\mathrm ds_i)\right)
       \left(\prod_{i=1}^r\alpha_{d-i+1,k}\delta_i^{-k}\right)
       \mathrm dA_{D_r}.
\end{equation}
Equality means equality after integration against every bounded Borel
function of $(A,\boldsymbol s)$.  Define the displayed product density
to be zero on degenerate tuples.  The measure on the right is supported
on tuples with $\dim D_i=i$ for all $1\le i\le r$.
\end{lemma}

\begin{proof}
We first take $P=\sigma$.  Fix $s_0,\ldots,s_{i-1}$ with
$\dim D_{i-1}=i-1$, and put $H=D_{i-1}^\perp$, so $\dim H=n_i=d-i+1$.
For a bounded continuous function $\varphi(A,s)$, consider
\begin{equation}\label{haar:small-ball-insertion}
 I_t=\int\mathrm dA_{D_{i-1}}\int\sigma(\mathrm ds)\,
       \varphi(A,s)
       \frac{\mathbf 1_{\{\lVert A(s-s_0)\rVert\le t\}}}{v_kt^k},
       \qquad t>0.
\end{equation}
We compute its limit as $t\downarrow0$ in both orders of integration.

Fix $A$ first.  Disintegration of $\sigma$ under $s\mapsto As$ rewrites
the inner integral as the average over the ball of radius $t$ about $As_0$
of
\[
 y\longmapsto p^0(y)\int\varphi(A,s)\,U_{A,y}(\mathrm ds).
\]
This function is continuous in the open unit ball.  Indeed, use the fiber
representation following \eqref{haar:posterior} with a common uniform
direction and apply bounded convergence.  Because $d>k+2$, the factor
$p^0$ tends to zero at the boundary, so the product extends continuously
by zero to $\mathbb R^k$.  The average therefore tends to
\begin{equation}\label{haar:first-order-limit}
 p^0(As_0)\int\varphi(A,s)\,U_{A,As_0}(\mathrm ds).
\end{equation}
Its absolute value, and that of every preceding average, is at most
$\lVert\varphi\rVert_\infty\alpha_{d,k}$.  Dominated convergence thus
permits the frame integration as well.

For the other order, fix $s$ and put
$\delta=\lVert\operatorname{proj}_H(s-s_0)\rVert$.  When $\delta>0$,
the vector $A(s-s_0)$ under $\mathrm dA_{D_{i-1}}$ has density
\begin{equation}\label{haar:insertion-density}
 z\longmapsto\delta^{-k}\alpha_{n_i,k}
       \left(1-\frac{\lVert z\rVert^2}{\delta^2}\right)^{(n_i-k-2)/2}
       \mathbf 1_{\{\lVert z\rVert<\delta\}}.
\end{equation}
This is \eqref{haar:sphere-density} in the current ambient space $H$,
not in the original $d$-dimensional space.  Since $n_i>k+2$, the density
is bounded by its value $\alpha_{n_i,k}\delta^{-k}$ at zero.

The conditional frame law also has a concrete limit.  Put
$e=\operatorname{proj}_H(s-s_0)/\delta$ and $w=z/\delta$.  If $C$ is a
Haar $k$-frame in $H\cap e^\perp$, then the conditional law of $A$ given
$A(s-s_0)=z$, for $\lVert w\rVert<1$, is the law of
\begin{equation}\label{haar:conditional-frame}
 A=w e^{\mathsf T}+(I_k-ww^{\mathsf T})^{1/2}C.
\end{equation}
Its rows are orthonormal and $Ae=w$.  Invariance under rotations fixing
$e$ identifies the remaining factor as Haar; conversely the displayed
formula describes every such frame.  The formula is continuous at $w=0$.
At zero the frame law is $\mathrm dA_{D_i}$, since
$H\cap e^\perp=D_i^\perp$.  It follows that the frame integral in
\eqref{haar:small-ball-insertion} tends, for this $s$, to
\begin{equation}\label{haar:second-order-limit}
 \alpha_{n_i,k}\delta^{-k}
       \int\varphi(A,s)\,\mathrm dA_{D_i}.
\end{equation}

The bound $\lVert\varphi\rVert_\infty\alpha_{n_i,k}\delta^{-k}$
dominates the frame averages for every $t$.  It is integrable in $s$.
Choose a $(k+1)$-dimensional subspace
$L\subset(D_{i-1}+\operatorname{span}\{s_0\})^\perp$; the available
dimension is at least $d-i\ge d-r>k+1$.  Then
$\delta\ge\lVert\operatorname{proj}_L s\rVert$.  The latter vector has
a bounded density by \eqref{haar:sphere-density}, because $d>k+3$.
In $k+1$ dimensions its inverse $k$th power is integrable at zero:
the radial integral is a constant times
$\int_0^1 u^k u^{-k}\,\mathrm du<\infty$.  This also shows that
$\delta>0$ for $\sigma$-almost every $s$.  Dominated convergence proves
the second limit after integration in $s$.

Equating the two integrated limits in \eqref{haar:first-order-limit} and
\eqref{haar:second-order-limit} gives the insertion identity.  Bounded
continuous functions determine finite Borel measures on the compact
frame--sphere space, so the identity also holds for bounded Borel tests.
The frame and fiber constructions above make the kernels measurable in
the preceding tuple, allowing successive integration.

Start with $\sigma(\mathrm ds_0)\mathrm dA$ and insert, for
$i=1,\ldots,r$, the kernels
$p^0(As_0)U_{A,As_0}(\mathrm ds_i)$.  Each insertion changes the frame
law from $\mathrm dA_{D_{i-1}}$ to $\mathrm dA_{D_i}$ and contributes
$\alpha_{d-i+1,k}\delta_i^{-k}$.  On the original side, disintegration
of $s_0$ contributes one more factor $p^0(y)$, so its total power is
$r+1=m$.  These measures are finite at every step because
$p^0\le\alpha_{d,k}$; in particular their final mass is
$\int p^0(y)^m\,\mathrm dy\le\alpha_{d,k}^{m-1}$.  The preceding
almost-sure positivity of every inserted height gives the full-rank
support.  This proves \eqref{haar:coarea-identity} for $\sigma$.

Finally multiply this equality of measures by the bounded Borel function
$\prod_{i=0}^r g(s_i)$.  Formula \eqref{haar:posterior} changes the
left side into the stated measure for $P$, while the right side becomes
the product of $P$ measures in \eqref{haar:coarea-identity}.  The bounded
tilt preserves finiteness.  No continuity of $g$ or of $p_A$ was used.
\end{proof}

We now pass from the finite measure to the actual posterior experiment.
Draw $(A,Y)$ with law $\mathrm dA\,p_A(y)\mathrm dy$, and, conditionally
on $(A,Y)$, draw $S_0,\ldots,S_r$ independently from $P_{A,Y}$.  Write
$Q$ for the law including the frame and tuple, and $Q_S$ for its tuple
marginal.  Each $(A,S_i)$ has law $\mathrm dA\,P(\mathrm ds)$, and all
the replicas satisfy $AS_i=Y$ almost surely.

Define the probability measure
$\Lambda(\mathrm dA,\mathrm d\boldsymbol s)
 =P^{\otimes m}(\mathrm d\boldsymbol s)\mathrm dA_{D_r}$.
The degenerate tuples have zero
$P^{\otimes m}$-measure, by the positivity proved during insertion; fix
any measurable convention there.  Multiply
\eqref{haar:coarea-identity} by $p_A(As_0)^{-r}$ on its positive set,
and by zero on its zero set.  This operation is justified first for the
minimum with a finite constant and then by monotone convergence.  On the
left it changes $p_A^m$ to $p_A$, since $m=r+1$ and the used posterior
is supported on $As_0=y$.  Consequently
\begin{equation}\label{haar:replica-density}
 \frac{\mathrm dQ}{\mathrm d\Lambda}(A,\boldsymbol s)
 =\left(\prod_{i=1}^r\alpha_{d-i+1,k}\delta_i^{-k}\right)
       p_A(As_0)^{-r}
\end{equation}
on the positive-density, full-rank set, with value zero elsewhere.

For a fresh Haar $\ell$-frame $B$, independent of this whole experiment,
put $Z_i=BS_i$ and let $P_B$ be the image of $P$ under $B$.  Define the
two total correlations
\begin{equation}\label{haar:correlations}
 J_S=D_{\mathrm{KL}}(Q_S\Vert P^{\otimes m}),\qquad
 J_Z=\mathbb E_B D_{\mathrm{KL}}
       (Q_{Z_0,\ldots,Z_r\mid B}\Vert P_B^{\otimes m}).
\end{equation}

\begin{lemma}[Finite upper correlation]\label{haar:upper-correlation}
The correlations satisfy $0\le J_Z\le J_S\le U<\infty$, where
\begin{align}
 U&=\sum_{i=1}^r\left[
       \log\alpha_{d-i+1,k}+k\mathbb E_Q\log(1/\delta_i)
       -\mathbb E_{A,S\sim P}\log p_A(AS)\right],
       \label{haar:upper-U}\\
 0\le U&\le Crd(\log K+\log d).\label{haar:crude-U}
\end{align}
Here $\mathbb E_{A,S\sim P}$ draws the Haar frame independently of a
vector with law $P$.  Every logarithm in \eqref{haar:upper-U} is
integrable.
\end{lemma}

\begin{proof}
The positive part of $\log p_A$ under $p_A$ is bounded uniformly in $A$
by $\log^+(K\alpha_{d,k})$.  Its negative part is also integrable
uniformly: the density is supported in a ball of finite volume, and
$-x\log x\le1/e$ for $0<x\le1$.  Comparison with uniform probability
on the unit ball now gives
\begin{equation}\label{haar:density-entropy}
 -\mathbb E_{A,S\sim P}\log p_A(AS)\le\log v_k.
\end{equation}

To control a height, condition on $(A,Y,S_0,\ldots,S_{i-1})$.  Put
$h=\sqrt{1-\lVert Y\rVert^2}$ and $n=d-k$.  Almost surely $h>0$.
The conditional residual direction of $S_i$ in $\ker A$ has density at
most
\[
 K'=\frac{Kp^0(Y)}{p_A(Y)}\ge1
\]
relative to uniform probability on $S^{n-1}$.  Choose a unit vector
$e\in\ker A$ perpendicular to the preceding residual directions,
including that of $S_0$.  There are at most $i\le r<n$ such directions;
a fixed Gram--Schmidt rule makes the choice measurable.  Then
$e\perp D_{i-1}$.  If $U_i$ is
the new unit residual direction, then
\[
 \delta_i\ge |\langle e,S_i-S_0\rangle|
             =h|\langle e,U_i\rangle|.
\]
The one-coordinate case of \eqref{haar:sphere-density} and
\eqref{haar:gamma-bounds} yield
\[
 \mathbb P\{|\langle e,U_i\rangle|\le t
       \mid A,Y,S_0,\ldots,S_{i-1}\}
 \le\min(1,CK'\sqrt d\,t).
\]
Integrating this tail with $t=e^{-u}$ shows that the conditional
expectation of $\log(1/|\langle e,U_i\rangle|)$ is at most
$1+\log(CK'\sqrt d)$.

For fixed $A$, the expected density tilt is
\[
 \mathbb E[\log K'\mid A]
  =\log K-D_{\mathrm{KL}}(p_A\Vert p^0)\le\log K.
\]
The divergence is finite: $p_A/p^0\le K$ controls its positive part,
and $-x\log x\le1/e$ controls its negative part under $p^0$.  Also the
individual marginal $(A,S_0)$ consists of an independent Haar frame and
a unit vector.  Thus $\lVert AS_0\rVert^2$ has beta parameters
$k/2,(d-k)/2$ for every fixed target, with both parameters positive.
The beta integral gives
\begin{equation}\label{haar:radius-moment}
 \mathbb E h^{-2}=\frac{d-2}{d-k-2}\le C,
 \qquad \mathbb E\log(1/h)\le\tfrac12\log\mathbb E h^{-2}\le C.
\end{equation}
Combining these estimates gives
$\mathbb E\log(1/\delta_i)\le C+\log K+C\log d$.
The negative part is bounded because $\delta_i\le2$.  Both signs of
every height logarithm are therefore integrable.

We may now take the expected logarithm of
\eqref{haar:replica-density}.  It is exactly the finite relative entropy
$D_{\mathrm{KL}}(Q\Vert\Lambda)=U$.
Forgetting $A$ by data processing yields $J_S\le U$.  Adding the
independent $B$ and then projecting coordinatewise yields $J_Z\le J_S$.
Finally \eqref{haar:gamma-bounds}, \eqref{haar:density-entropy}, and
the height bound give \eqref{haar:crude-U}.
\end{proof}

The inverse-height terms in the upper bound can grow with the density
bound $K$.  We will keep those terms until the projected correlation
has supplied matching lower bounds.  This order is what permits averaging
over rare states at the end.

\subsection{The correlation retained by the auxiliary projection}

We will bound the probability that an independent projected signal lies
close to the affine span of preceding projected replicas.  Averaging over
the side frame turns this tube probability into ball averages of $p_A$
around the actual label $AS_0$.  Their radii vary with the simplex height;
the next lemma controls the logarithmic cost of taking a supremum over
those radii.

For a probability density $f$ on $\mathbb R^k$, define its centered
maximal average by
\[
 \mathcal M f(u)=\sup_{t>0}\frac{1}{v_kt^k}
        \int_{\lVert u'-u\rVert\le t} f(u')\,\mathrm du'.
\]
For fixed $f\in L^1$, the ball integral is continuous in its center and
positive radius; the supremum may therefore be restricted to positive
rational radii.  In particular $\mathcal M f$ is measurable.  For the
jointly measurable family $p_A$, the same rational supremum and Tonelli's
theorem give joint measurability in $(A,u)$.

For historical background on one-dimensional maximal averages, see Hardy
and Littlewood~\cite[Section~III]{HardyLittlewood1930}.  Stein and
Str\"omberg~\cite[p.~259]{SteinStromberg1983} discuss the exponential
dimension cost of customary covering proofs.  We prove the centered weak
estimate needed here, including its constant.

\begin{lemma}[Logarithmic maximal estimate]\label{haar:maximal}
Suppose $f$ is a probability density supported in the unit ball of
$\mathbb R^k$ and $f\le K\alpha_{d,k}$.  Then
\begin{equation}\label{haar:maximal-log}
 \int f(u)\log\frac{\mathcal M f(u)}{f(u)}\,\mathrm du
 \le Cd+\log(1+Cd+\log K).
\end{equation}
Both logarithms are integrable under $f(u)\mathrm du$.
The integrand is taken as zero on $\{f=0\}$.
\end{lemma}

\begin{proof}
For every $\lambda>0$,
\begin{equation}\label{haar:weak-maximal}
 \bigl|\{u:\mathcal M f(u)>\lambda\}\bigr|
       \le\frac{3^k}{\lambda}.
\end{equation}
To prove this, take a compact subset of the displayed superlevel set.
Open and closed balls have the same $f$-mass, since their boundaries have
Lebesgue measure zero.
At each point choose an open ball centered there whose $f$-mass exceeds
$\lambda$ times its volume.  The strict inequality persists under a
sufficiently small enlargement, so these balls give an open cover.
Take a finite subcover and greedily keep disjoint balls in decreasing
order of radius.  Every discarded ball meets a retained ball of at least
its radius, and is contained in that ball's threefold dilation.  Hence
the compact set has volume at most $3^k$ times the sum of the retained
volumes.  The retained balls are disjoint and have total $f$-mass at most
one, so their volumes sum to less than $1/\lambda$.  Inner regularity
proves \eqref{haar:weak-maximal}.

Set $b=3^k$ and $L_f=K\alpha_{d,k}$.  Since $\mathcal M f\le L_f$,
layer integration over the unit ball gives
\begin{align*}
 \int_{\lVert u\rVert\le1}\mathcal M f(u)\,\mathrm du
 &\le\int_0^{L_f}\min(v_k,b/\lambda)\,\mathrm d\lambda\\
 &\le b\bigl(1+\log^+(L_fv_k/b)\bigr).
\end{align*}
On the support of $f$, a radius-two ball contains the whole unit ball,
so $\mathcal M f\ge(v_k2^k)^{-1}$ there.  Its logarithm is bounded on
that support, while $f\log f$ is integrable by the bounded-density
argument in Lemma~\ref{haar:upper-correlation}.  Jensen's inequality
under the probability density $f$ now yields
\[
 \int f\log(\mathcal M f/f)
 \le\log\int_{\{f>0\}}\mathcal M f
 \le k\log3+\log\bigl(1+\log^+(K\alpha_{d,k}v_k/3^k)\bigr).
\]
Use \eqref{haar:gamma-bounds} and $d/k\le32$ to obtain
\eqref{haar:maximal-log}.
\end{proof}

The next test uses only the projected points, including when their
conditional distance distribution has atoms.

\begin{lemma}[Lower projected correlation]\label{haar:lower-correlation}
For the replica law above and every $0<a<1$,
\begin{equation}\label{haar:lower-Z}
 J_Z\ge a\sum_{i=1}^r\left[
       k\mathbb E_Q\log(1/\delta_i)-\log(C_1^kv_k)
       -\mathbb E_{A,S\sim P}\log\mathcal M p_A(AS)\right]
       +r\log(1-a),
\end{equation}
where $C_1$ is an absolute constant, independent of $P,K,d,a$.
\end{lemma}

\begin{proof}
Given $(B,Z_0,\ldots,Z_{i-1})$, define for $z\in\mathbb R^\ell$
\[
 R_i(z)=\operatorname{dist}\bigl(z-Z_0,
             \operatorname{span}\{Z_j-Z_0:1\le j<i\}\bigr).
\]
Let $F_i(t)=P\{s:R_i(Bs)\le t\}$ be its distribution function under
an independent $s\sim P$.
The distance, and hence this conditional distribution function, is
measurable in the projected history: the distance is the infimum of
the corresponding norms over rational coefficient vectors.  Let
$\rho_i=R_i(Z_i)$ be the actual replica distance.

If a random variable $R$ has distribution function $F$, then
\begin{equation}\label{haar:superuniform}
 \mathbb P\{F(R)\le u\}\le u\qquad(0\le u\le1).
\end{equation}
This includes distributions with atoms.  In fact, the set of real $x$
with $F(x)\le u$ is an initial interval, whose probability is the
supremum of $F$ on that interval and is at most $u$.  Tail integration
of \eqref{haar:superuniform} gives, for $0<a<1$,
\begin{equation}\label{haar:cdf-moment}
 \mathbb E F(R)^{-a}
 \le1+\int_1^\infty t^{-1/a}\,\mathrm dt=\frac1{1-a}.
\end{equation}
In particular $F(R)>0$ almost surely under its own law.

Apply the relative-entropy variational inequality to the test
$z\mapsto F_i(R_i(z))^{-a}$ under the conditional one-vector reference
law $P_B$.
For an unbounded test, use its minimum with a finite constant and pass
monotonically to the limit; \eqref{haar:cdf-moment} bounds its reference
expectation.  The chain rule for $J_Z$ has one term for each conditional
law of $Z_i$ given $(B,Z_0,\ldots,Z_{i-1})$, and its $i=0$ term is zero.
It follows that
\begin{equation}\label{haar:cdf-test}
 J_Z\ge a\sum_{i=1}^r\mathbb E_{Q,B}\log\frac1{F_i(\rho_i)}
             +r\log(1-a).
\end{equation}
The logarithmic expectations are finite: the variational bound after
averaging controls their nonnegative parts by the finite correlation
$J_Z\le J_S$ from Lemma~\ref{haar:upper-correlation}.

We next estimate $F_i(\delta_i)$ while holding the entire replica
outcome $(A,S_0,\ldots,S_r)$ fixed.  Put $D=D_{i-1}$ and
$E=\operatorname{row}(B)$.  The isometry $B^{\mathsf T}:\mathbb R^\ell
\to E$ gives the exact identity
\begin{equation}\label{haar:intersection}
 \operatorname{dist}(Bu,BD)
       =\lVert\operatorname{proj}_{E\cap D^\perp}u\rVert.
\end{equation}
Indeed, the orthogonal complement within $E$ of
$\operatorname{proj}_E D$ is precisely $E\cap D^\perp$.
The projection of the fixed $(i-1)$-space $D$ into a Haar $\ell$-space
$E$ has full rank almost surely: generate $E$ with independent standard
Gaussians and restrict them to an orthonormal basis of $D$, obtaining an
$\ell\times(i-1)$ Gaussian matrix of full column rank.  Thus
$\dim(E\cap D^\perp)=q_i=\ell-i+1$.  Rotations preserving $D$ show
that this intersection is a uniform $q_i$-space in the $n_i$-space
$D^\perp$.

For fixed $u$, let $v=\operatorname{proj}_{D^\perp}u$.  Spherical
projection in that $n_i$-space, followed by
\eqref{haar:gamma-bounds}, gives
\begin{equation}\label{haar:projected-small-ball}
 \mathbb P_B\{\operatorname{dist}(Bu,BD)\le t\}
 \le\min\left(1,\left(\frac{C_0t}{\lVert v\rVert}\right)^{q_i}\right),
\end{equation}
with value one if $v=0$.  The constant $C_0$ is absolute because
$n_i/q_i\le16/3$; the density formula is applicable because
$n_i-q_i=d-\ell>2$.  In the replica experiment $A$ annihilates $D$.
Therefore for $u=s-S_0$,
$\lVert v\rVert\ge\lVert A(s-S_0)\rVert$.  Also
$\rho_i\le\delta_i$, since $B$ is a contraction and the distance
minimization permits all vectors in $D$.

Set $u_0=AS_0$ and $q=q_i$.  The identity
\[
 \min(1,(b/x)^q)=\int_1^\infty qt^{-q-1}
                 \mathbf 1_{\{x\le bt\}}\,\mathrm dt
       \qquad(b>0,\ x\ge0)
\]
also holds at $x=0$.  Integrating \eqref{haar:projected-small-ball}
against $s\sim P$ and using Tonelli's theorem yields
\begin{align}
 \mathbb E_B F_i(\delta_i)
 &\le\int_1^\infty qt^{-q-1}
       P\{\lVert As-u_0\rVert\le C_0\delta_i t\}\,\mathrm dt\notag\\
 &\le v_k(C_0\delta_i)^k\mathcal M p_A(u_0)
             \int_1^\infty qt^{k-q-1}\,\mathrm dt\notag\\
 &=\frac{q_i}{q_i-k}v_k(C_0\delta_i)^k\mathcal M p_A(u_0)
 \le C_1^kv_k\delta_i^k\mathcal M p_A(u_0).
 \label{haar:projected-maximal}
\end{align}
The integral converges because $q_i>k$, and the last step uses the exact
bound $q_i/(q_i-k)\le3/2$.  One absolute choice of $C_1$ works for every
$k\ge1$.

Now $F_i(\rho_i)\le F_i(\delta_i)$.  Jensen's inequality for $-\log$
and \eqref{haar:projected-maximal}, first for each fixed replica outcome,
give
\[
 \mathbb E_{Q,B}\log\frac1{F_i(\rho_i)}
 \ge k\mathbb E_Q\log(1/\delta_i)-\log(C_1^kv_k)
       -\mathbb E_{A,S\sim P}\log\mathcal M p_A(AS).
\]
Here the last expectation uses the individual marginal $(A,S_0)$.
The right side is finite by Lemmas~\ref{haar:upper-correlation} and
\ref{haar:maximal}; Jensen may also be read in the extended sense before
averaging.  Substitution into \eqref{haar:cdf-test} proves
\eqref{haar:lower-Z}.
\end{proof}

\subsection{Cancellation and the block information bound}

The preceding estimates can now be subtracted.  Retaining every term of
$U$ gives, for every $0<a<1$,
\begin{align}
 J_S-J_Z\le{}&(1-a)U-r\log(1-a)\notag\\
 &+a\sum_{i=1}^r\left[
       \log(\alpha_{d-i+1,k}C_1^kv_k)
       +\mathbb E_{A,S\sim P}
          \log\frac{\mathcal M p_A(AS)}{p_A(AS)}\right].
 \label{haar:cancellation}
\end{align}
Thus only the fraction $a$ of the height contribution has been canceled;
the residual $(1-a)U$ remains explicit.  Take
\begin{equation}\label{haar:a-choice}
 a=1-d^{-3}.
\end{equation}
Using \eqref{haar:crude-U}, \eqref{haar:gamma-bounds}, and
Lemma~\ref{haar:maximal} in \eqref{haar:cancellation} yields
\begin{equation}\label{haar:fixed-prior-budget}
 \frac{J_S-J_Z}{m}
 \le Cd+\log(1+Cd+\log K)
       +C\frac{\log K+\log d}{d^2}.
\end{equation}
In detail, $(1-a)U/m\le C(\log K+\log d)/d^2$ because $r/m\le1$,
and $-r\log(1-a)/m\le3\log d$.  Also
$\log(\alpha_{n_i,k}C_1^kv_k)\le Cd$ by
\eqref{haar:gamma-bounds} and $n_i/k\le32$.  These are all the terms in
\eqref{haar:cancellation}; the $3\log d$ term is absorbed into $Cd$.

\begin{proposition}[Haar block information bound]\label{haar:block}
Let $S\sim\sigma$, let $V$ be a finite-valued state, and let a block
consist of $k=\lfloor d/16\rfloor$ fresh independent standard Gaussian
designs, independent of $(S,V)$ and the past, with their exact labels.
Suppose a finite-valued state $W$ is
generated by a measurable randomized kernel from $V$ and the full block,
with no other retained information from the past.  Let $B$ be a Haar
$\ell=\lfloor d/4\rfloor$ frame drawn independently of the entire joint
experiment producing $S,V$, the block, and $W$.  Then
\begin{align}
 I(S;W\mid V,B,BS)\le{}&\frac{H(W\mid V)}{m}+Cd
       +\log(1+Cd+H(V))\notag\\
 &+C\frac{H(V)+\log d}{d^2},
 \qquad m=\lfloor d/16\rfloor+1.
 \label{haar:state-budget}
\end{align}
The constant $C$ is absolute.  If every state has logarithmic width at
most $H_*=o(d^2)$, the left side is at most $Cd$ for all sufficiently
large $d$.
\end{proposition}

\begin{proof}
Fix a state value $v$ with probability $p_v>0$.  Its conditional signal
law has the bounded density
\[
 g_v(s)=\frac{\mathbb P\{V=v\mid S=s\}}{p_v}\le K_v:=\frac1{p_v}
\]
relative to $\sigma$.  The new Gaussian design matrix is independent of
$(S,V)$.  Its row Gram--Schmidt factorization is $RA$ with an invertible
triangular $R$ and a Haar $k$-frame $A$, almost surely.  The exact labels
are $R(AS)$, so, conditional on the design, observing them is equivalent
to observing $Y=AS$.  Because the whole design is independent of $S$
conditional on $V=v$, its coefficients $R$ add no information about $S$
once $(A,Y)$ is specified.

Generate the replicas conditionally independently from this posterior,
and independently of $W$ given $v$ and the complete block data.  A single
replica then has the same joint law with the block data and $W$ as the
original signal: both use the posterior given the block, and $W$ has no
further dependence on the signal once those data and $v$ are fixed.
In particular each pair $(S_i,W)$ has the same conditional law as
$(S,W)$, the auxiliary $B$ is independent of the entire tuple and $W$,
and the original pair used on the left is also coupled independently
with $B$.  The common map $(B,s)\mapsto Bs$ is continuous, and the
unprojected total correlation is finite by
Lemma~\ref{haar:upper-correlation}.  These are precisely the application
hypotheses of Lemma~\ref{lem:replica-information}.  Applied in the
conditional experiment $V=v$, it gives
\[
 mI(S;W\mid V=v,B,BS)
 \le H(W\mid V=v)+J_S(v)-J_Z(v).
\]
No conditional independence of the replicas given $W$ is asserted or
needed.

Insert \eqref{haar:fixed-prior-budget} with $K=K_v$ and average over $v$.
The identity
\[
 \sum_v p_v\log K_v=\sum_vp_v\log(1/p_v)=H(V)
\]
accounts exactly for rare state values.  Jensen's inequality bounds the
average of $\log(1+Cd+\log K_v)$ by $\log(1+Cd+H(V))$.  This proves
\eqref{haar:state-budget}.  Under the width hypothesis,
$H(W\mid V)/m=o(d)$, the last term is $o(1)$, and the logarithmic term
is $O(\log d)$ for all sufficiently large $d$.  The stated $Cd$ bound
follows.
\end{proof}

\subsection{The terminal fiber and the streaming conclusion}

The block comparison is complete.  We apply the residual-sphere endpoint
of Section~\ref{sec:model-endpoint} with the present Haar side frame,
then compare the resulting information with the number of blocks.

\begin{lemma}[The terminal Haar fiber]\label{haar:terminal}
Let $S\sim\sigma$, let $V_f$ be finite-valued, and let the unit-valued
output $\widehat S$ be generated from $V_f$ by a measurable kernel using
fresh randomness, conditionally independently of $S$.  Let $B$ be a Haar
$\ell=\lfloor d/4\rfloor$ frame independent of $(S,V_f,\widehat S)$.
If
$\mathbb P\{\arccos\langle\widehat S,S\rangle\le\epsilon\}\ge2/3$
for $0<\epsilon\le1/10$, then
\begin{equation}\label{haar:terminal-information}
 I(S;V_f\mid B,BS)
 \ge\frac{d-\ell-1}{3}\log\frac1{4\epsilon}-\log2.
\end{equation}
\end{lemma}

\begin{proof}
Apply Lemma~\ref{lem:residual-information} with $G=B$, $V=V_f$,
$b=\ell$, and $r_0=1/2$.  Its hypotheses hold because $B$ is independent of
$(S,V_f,\widehat S)$.  Explicitly, the Markov
bound for the residual radius $h$ is
$\mathbb P\{h<1/2\}\le4\ell/(3d)\le1/3$.  Success together with
$h\ge1/2$ therefore has probability at least $1/3$, while the conditional
prior cap probability is at most $(4\epsilon)^{d-\ell-1}$.  These are exactly
the constants in \eqref{haar:terminal-information}.
\end{proof}

\begin{proposition}[Streaming consequence]\label{haar:streaming}
For each dimension $d$, let $S\sim\sigma$ and consider a learner in
Definition~\ref{def:streaming-model}, with deterministic horizon $T=T(d)$
and nonnegative integer memory $M=M(d)$.  If $M(d)=o(d^2)$,
$0<\epsilon(d)\le1/10$, and its angular success probability averaged over
the signal, samples, and learner randomness is at least $2/3$, then
\[
 T(d)=\Omega\bigl(d\log(1/\epsilon(d))\bigr).
\]
The implicit constant is absolute; the dimension threshold may depend on
the memory sequence.
\end{proposition}

\begin{proof}
Apply the fixed-shared-seed part of Lemma~\ref{lem:padding} to the learner
in Definition~\ref{def:streaming-model}.  The resulting Borel experiment
has data-independent initialization and uniform-prior success at least
$2/3$, with fresh transition and output randomness retained.  This is
the success threshold in Lemma~\ref{haar:terminal}.

First, success forces $T=\Omega(d)$ without any memory restriction.
Suppose $T\le d/8$, and grant an estimator all $T$ raw pairs, filling
in ignored samples after any early stopping.  Lemma~\ref{lem:residual-sphere}
with $b=T$ gives a uniform residual sphere in kernel dimension $d-T$;
\eqref{eq:residual-markov} bounds its probability of radius below $1/2$
by $4T/(3d)$.  At $T=0$ the radius is one deterministically.  On radius
at least $1/2$, the cap argument of Lemma~\ref{lem:raw-sample-bound}
bounds conditional success by $(4\epsilon)^{d-T-1}$.  Thus
\begin{equation}\label{haar:raw-sample-bound}
 \mathbb P\{\text{angular error}\le\epsilon\}
 \le\frac{4T}{3d}+(4\epsilon)^{d-T-1}
       \qquad(T\le d/8).
\end{equation}
The output has no further dependence on $S$ conditional on all the raw
data, so the bound also covers randomized estimators.  Its right side is
at most $1/6+(0.4)^{7d/8-1}<2/3$ for all sufficiently large $d$.
Hence $T>d/8$ on those dimensions.

Put $L=\log(1/\epsilon)\ge\log10$.  If $T\ge dL$, the conclusion already
holds.  On the remaining branch $T<dL$, apply the output-capacity bound
of Lemma~\ref{lem:early-clock} to this uniform-prior experiment with
$p_0=2/3$.  It gives
\begin{equation}\label{haar:output-count}
 (d-1)L\le M\log2+\log(T+1)+\log(3/2).
\end{equation}
For large $d$, $\log(T+1)\le\log(dL+1)\le dL/2$: here $dL\ge d\log10$
and $\log(x+1)\le x/2$ for $x\ge3$.  Thus
\eqref{haar:output-count} implies
\begin{equation}\label{haar:short-clock}
 dL\le C(M+1),\qquad L=o(d),\qquad T=o(d^2).
\end{equation}
For example, after subtracting $dL/2$, the left side is
$(d/2-1)L\ge dL/4$ for $d\ge4$.  The second and third conclusions follow
from $M=o(d^2)$ and $T<dL$.

Use Lemma~\ref{lem:padding} to retain the stopping index, postpone the
fresh randomized output, and append ignored rows through the last block.
Its state count gives every padded state entropy at most
\begin{equation}\label{haar:padded-width}
 H_*=M\log2+\log(T+2)
 \le M\log2+\log(C(M+1)+2)=o(d^2).
\end{equation}
This count restricts the branch where the block argument is needed; it
does not replace that argument in the representable-precision regime.

Use a single independent Haar $\ell$-frame $B$ for the entire analysis,
put $Z=BS$, and divide the padded stream into
$N=\lceil T/k\rceil$ blocks.  For consecutive states $V,W$, the chain
rule and data processing give
\[
 I(S;W\mid B,Z)
 \le I(S;V,W\mid B,Z)
 =I(S;V\mid B,Z)+I(S;W\mid V,B,Z).
\]
The initial state is independent of $(S,B,Z)$.  Proposition~\ref{haar:block}
and \eqref{haar:padded-width} therefore imply
\begin{equation}\label{haar:information-clock}
 I(S;V_f\mid B,Z)\le Cd\lceil T/k\rceil.
\end{equation}
The terminal lower bound \eqref{haar:terminal-information} is at least
$c dL$ for all sufficiently large $d$: indeed
$d-\ell-1\ge d/2$ and
\[
 \log\frac1{4\epsilon}
 \ge\left(1-\frac{\log4}{\log10}\right)L,
\]
while the subtractive $\log2$ is absorbed because $dL\to\infty$.
Finally $k\ge d/32$ and the already proved $T>d/8$ give
\[
 d\lceil T/k\rceil\le32T+d\le40T.
\]
Comparison with \eqref{haar:information-clock} proves the proposition.
\end{proof}

A learner that succeeds with probability at least $2/3$ for every fixed
signal also has uniform-prior success at least $2/3$ by integration.
Thus Proposition~\ref{haar:streaming} separately gives the corresponding
worst-signal lower bound, with the same exact observations and finite-state
model.
